\chapter{Acuerdo de Confidencialidad del Social Guide}
\label{anexo:confidencialidad}

\section[Declaración de Confidencialidad]{Declaración de Confidencialidad y\\Compromiso Ético}

Yo, \underline{\hspace{6cm}} (nombre completo del Social Guide), en mi rol de \textbf{Social Guide} del equipo \underline{\hspace{5cm}} (nombre del equipo), me comprometo formalmente a cumplir con los principios éticos y salvaguardas de confidencialidad descritos a continuación, en el marco de la implementación de SocialScrum.

\subsection{Confidencialidad absoluta de información individual}

\begin{itemize}
    \item \textbf{No compartiré información que permita identificar a personas específicas} fuera del equipo, incluyendo, pero sin limitarse a:
          \begin{itemize}
              \item Comentarios realizados durante eventos sociales (Social Daily Standup, Social Sprint Retrospective, sesiones de mediación).
              \item Conflictos interpersonales descritos por miembros del equipo.
              \item Información relacionada con el estado emocional, \textit{burnout}, desmotivación o vulnerabilidades personales.
              \item Cualquier dato que pueda revelar directa o indirectamente la identidad de un miembro del equipo.
          \end{itemize}

    \item \textbf{Solo reportaré información agregada y anonimizada}, como:
          \begin{itemize}
              \item El Health Score promedio del equipo (escala 1–10).
              \item \gls{community-smells} identificados a nivel de equipo, sin atribuir responsabilidades individuales.
              \item Tendencias de salud social a lo largo del tiempo.
              \item Necesidades de intervención organizacional sin mencionar casos particulares.
          \end{itemize}
\end{itemize}

\subsection{Estructura de reporte y rendición de cuentas}

\textbf{Regla fundamental}: El Social Guide nunca reporta a Recursos Humanos (RR. HH.). Reporta al CTO/VP Engineering a través del Scrum Master.

\subsubsection{Tipos de reportes}

\begingroup
\scriptsize
\renewcommand{\arraystretch}{1.4}
\setlength{\tabcolsep}{4pt}
\setlength{\LTleft}{0pt}
\setlength{\LTright}{0pt}

\begin{longtable}[H]{
    >{\RaggedRight\hspace{0pt}}p{2.8cm}
    >{\Centering\hspace{0pt}}p{2.5cm}
    >{\RaggedRight\hspace{0pt}}p{6.9cm}
    >{\Centering\hspace{0pt}}p{2.5cm}
    }

    \caption{Reportes del Social Guide.} \label{tab:reportes-social-guide}                                                                                                     \\
    \toprule
    \textbf{Tipo}                                                                                              & \textbf{Frecuencia} & \textbf{Contenido} & \textbf{Audiencia} \\
    \midrule
    \endfirsthead

    \multicolumn{4}{c}{{\bfseries \tablename\ \thetable{} -- Continuación}}                                                                                                    \\
    \toprule
    \textbf{Tipo}                                                                                              & \textbf{Frecuencia} & \textbf{Contenido} & \textbf{Audiencia} \\
    \midrule
    \endhead

    \midrule
    \multicolumn{4}{r}{\textit{Continúa en la siguiente página...}}                                                                                                            \\
    \endfoot

    \bottomrule
    \endlastfoot

    Interno al equipo                                                                                          &
    Cada Sprint                                                                                                &
    Health Score del equipo, \gls{community-smells} activos y acciones sociales completadas durante el Sprint. &
    Equipo                                                                                                                                                                     \\
    \midrule

    Ejecutivo                                                                                                  &
    Mensual                                                                                                    &
    Health Score promedio, \gls{community-smells} críticos y solicitudes de recursos o apoyo organizacional.   &
    CTO, Scrum Master                                                                                                                                                          \\
\end{longtable}
\endgroup

Todos los reportes son \textbf{agregados y anónimos}. No identifican a personas.

\subsubsection{Protocolo si RR. HH. solicita información}

\begin{enumerate}
    \item El Social Guide rechaza la solicitud inmediatamente sin proporcionar información.
    \item Informa al Scrum Master y al CTO de la solicitud.
    \item El CTO evalúa si existe justificación legal (ej. investigación formal de acoso).
    \item Si existe justificación legal: se proporciona información estrictamente relevante, con consentimiento documentado.
    \item Si no existe justificación legal: el CTO responde que la información está protegida por políticas de confidencialidad.
\end{enumerate}

\subsubsection{Rendición de cuentas}

\begin{itemize}
    \item \textbf{Auditoría semestral}: Un auditor externo entrevista al equipo, revisa los reportes y emite un informe público.
    \item \textbf{Feedback continuo}: Encuesta anónima cada Sprint (evaluando confidencialidad, seguridad, neutralidad). Si la aprobación es inferior al 60~\% durante 2 Sprints consecutivos, se procede a una revisión formal.
\end{itemize}

\subsection{Dedicación del rol}

\begingroup
\scriptsize
\renewcommand{\arraystretch}{1.4}
\setlength{\tabcolsep}{4pt}
\setlength{\LTleft}{0pt}
\setlength{\LTright}{0pt}

\begin{longtable}[H]{
    >{\RaggedRight\hspace{0pt}}p{7.2cm}
    >{\Centering\hspace{0pt}}p{3cm}
    >{\RaggedRight\hspace{0pt}}p{5cm}
    }

    \caption{Dedicación del Social Guide según contexto.} \label{tab:dedicacion-social-guide} \\
    \toprule
    \textbf{Contexto}                                 & \textbf{Dedicación} & \textbf{Modelo} \\
    \midrule
    \endfirsthead

    \multicolumn{3}{c}{{\bfseries \tablename\ \thetable{} -- Continuación}}                   \\
    \toprule
    \textbf{Contexto}                                 & \textbf{Dedicación} & \textbf{Modelo} \\
    \midrule
    \endhead

    \midrule
    \multicolumn{3}{r}{\textit{Continúa en la siguiente página...}}                           \\
    \endfoot

    \bottomrule
    \endlastfoot

    Equipo de 5--7 personas con Health Score $\geq 7$ &
    10--20~\%                                         &
    Rol rotativo                                                                              \\
    \midrule

    Equipo de 5--7 personas con Health Score $< 7$    &
    30--40~\%                                         &
    Rol parcial                                                                               \\
    \midrule

    Equipo de 8--15 personas                          &
    30--70~\%                                         &
    Rol parcial o casi completo                                                               \\
    \midrule

    Equipo de 16 o más personas                       &
    80--100~\%                                        &
    Tiempo completo                                                                           \\
    \midrule

    Múltiples equipos                                 &
    100~\%                                            &
    Dedicado a 2--3 equipos                                                                   \\
\end{longtable}
\endgroup

\subsection{Separación total entre salud social y evaluación de desempeño}

\begin{itemize}
    \item \textbf{No utilizaré} la información que recopile como Social Guide para:
          \begin{itemize}
              \item Evaluaciones de desempeño individuales.
              \item Decisiones de contratación, promoción, reasignación o despido.
              \item Justificar sanciones disciplinarias.
              \item Influir en decisiones salariales o de compensación.
          \end{itemize}

    \item \textbf{Rechazaré cualquier solicitud de información individual}, incluso si proviene de:
          \begin{itemize}
              \item Recursos Humanos (RR. HH.), salvo los casos excepcionales previstos en la siguiente sección.
              \item Gerencia o liderazgo técnico, cuando dicha solicitud viole la confidencialidad acordada.
              \item Entidades externas a la organización.
          \end{itemize}
\end{itemize}

\subsection{Excepciones legalmente justificadas}

Reconozco que existen situaciones excepcionales que requieren suspender la confidencialidad:

\begin{itemize}
    \item \textbf{Riesgo inminente de daño}: si un miembro expresa intención de autolesión, daño a terceros o conducta ilegal, debo reportarlo a las autoridades correspondientes (RR. HH., seguridad, etc.).

    \item \textbf{Investigaciones formales por acoso o discriminación}: en estos casos:
          \begin{itemize}
              \item Solo compartiré información con el consentimiento explícito de la(s) persona(s) afectada(s), cuando sea posible.
              \item Me limitaré estrictamente a los hechos pertinentes.
              \item Documentaré la excepción y notificaré al equipo siempre que sea legalmente permitido.
          \end{itemize}
\end{itemize}

Siempre que la ley lo permita, \textbf{informaré al equipo} cuando deba suspender la confidencialidad y explicaré las razones.

\subsection{Seguridad psicológica, auditoría y mejora continua}

Me comprometo a proteger la seguridad psicológica del equipo, actuando con integridad y facilitando sin imponer. El equipo tiene derecho a revisar cualquier reporte antes de su difusión, auditar mi cumplimiento en cualquier momento y solicitar mi remoción inmediata si detecta una violación. Me comprometo a participar en auditorías semestrales independientes, recibir retroalimentación del equipo y mantener formación continua en ética y facilitación. Los fundamentos detallados de estas salvaguardas se describen en el Capítulo~\ref{chap:capitulo3}, Sección~\ref{sec:salvaguardas_eticas}.

\subsection{Consecuencias de la violación de este acuerdo}

Reconozco que, si incumplo alguno de los principios establecidos:

\begin{itemize}
    \item Seré \textbf{removido de inmediato} del rol de Social Guide.
    \item \textbf{No podré ejercer este rol nuevamente} dentro de la organización.
    \item La organización ofrecerá una \textbf{disculpa formal} al equipo.
    \item Se realizará una \textbf{investigación formal} y se tomarán medidas correctivas si hubo daño comprobado.
\end{itemize}

\vspace{1cm}

\noindent
\textbf{Firma del Social Guide:} \underline{\hspace{6cm}} \\[0.5cm]
\textbf{Nombre completo:} \underline{\hspace{6cm}} \\[0.5cm]
\textbf{Fecha:} \underline{\hspace{4cm}} \\[1cm]

\noindent
\textbf{Testigo (Scrum Master):} \underline{\hspace{6cm}} \\[0.5cm]
\textbf{Nombre completo:} \underline{\hspace{6cm}} \\[0.5cm]
\textbf{Fecha:} \underline{\hspace{4cm}} \\[1cm]

\noindent
\textbf{Testigo (CTO/VP Ingeniería):} \underline{\hspace{6cm}} \\[0.5cm]
\textbf{Nombre completo:} \underline{\hspace{6cm}} \\[0.5cm]
\textbf{Fecha:} \underline{\hspace{4cm}} \\